\section*{Bab \ref{ch:b1:photosynthesis} --- Fotosintesis dan Autotrofi}

\begin{solution}{exo:b1:photosynthesis:1}
Membran selubung luar dan dalamnya, serta membran \omterm{def:b1:photosynthesis:chloroplast}{tilakoid} di dalamnya.
\omterm{prop:b1:photosynthesis:lightreactions}{Reaksi terang}: dalam membran \omterm{def:b1:photosynthesis:chloroplast}{tilakoid} (proton dipompa ke dalam
lumennya). \omterm{def:b1:photosynthesis:fixation}{Daur Calvin}: dalam \omterm{def:b1:photosynthesis:chloroplast}{stromanya}.
\end{solution}

\begin{solution}{exo:b1:photosynthesis:2}
$2\,\mathrm{H_2O} + 2\,\mathrm{NADP^+} + 3\,\mathrm{ADP} + 3\,\mathrm{P_i}
\to \mathrm{O_2} + 2\,\mathrm{NADPH} + 2\,\mathrm{H^+} + 3\,\mathrm{ATP}$
dengan delapan foton. Oksigennya datang dari air: alga yang diberi
$\mathrm{H_2^{18}O}$ melepas $\mathrm{^{18}O_2}$; yang diberi
$\mathrm{C^{18}O_2}$ tidak (Ruben dan Kamen).
\end{solution}

\begin{solution}{exo:b1:photosynthesis:3}
\omterm{def:b1:photosynthesis:fixation}{Fiksasi} (\omterm{def:b1:photosynthesis:fixation}{RuBisCO}), reduksi (menjadi G3P), pembentukan kembali (RuBP).
Tiap $\mathrm{CO_2}$: 3 \omterm{def:b1:water-small-molecules:atp}{ATP} dan 2 NADPH.
\end{solution}

\begin{solution}{exo:b1:photosynthesis:4}
Dekat \qty{430}{nm} (biru) dan \qty{660}{nm} (merah). Di sekitar
\qty{500}{nm} \omterm{def:b1:photosynthesis:pigments}{karotenoidnya} menyerap lalu meneruskan energinya ke
\omterm{def:b1:photosynthesis:pigments}{klorofil}, sehingga oksigen dilepas meskipun \omterm{def:b1:photosynthesis:pigments}{klorofilnya} sendiri sedikit
menyerap di sana.
\end{solution}

\begin{solution}{exo:b1:photosynthesis:5}
Cahaya merah jauh (\qty{700}{nm}) sendirian sedikit menggerakkan
fotosintesis; ditambahkan pada cahaya \qty{680}{nm} ia memberikan lebih
daripada jumlah kedua lajunya. Dua sistem pigmen dengan puncak serapan
yang berbeda pasti \omterm{def:b1:proteins:allostery}{bekerja sama} berderet, masing-masing memerlukan
fotonnya sendiri: \omterm{def:b1:photosynthesis:zscheme}{fotosistem} I (P700) dan \omterm{def:b1:photosynthesis:zscheme}{fotosistem} II (P680).
\end{solution}

\begin{solution}{exo:b1:photosynthesis:6}
$0.18\times 96\,500 = \qty{17.4}{kJ/mol}$ proton; $50/17.4 = 2.9$
proton sekurang-kurangnya; sintasenya memakai empat, yakni efisiensi
\qty{72}{\%}.
\end{solution}

\begin{solution}{exo:b1:photosynthesis:7}
Ia membuktikan bahwa sebuah gradien proton melintasi membran \omterm{def:b1:photosynthesis:chloroplast}{tilakoid}
sudah cukup untuk menggerakkan sintesis \omterm{def:b1:water-small-molecules:atp}{ATP}, tanpa cahaya dan tanpa
pengangkutan elektron: gandengannya lewat gradiennya. Ia sendirian
tidak membuktikan bahwa \omterm{prop:b1:photosynthesis:lightreactions}{reaksi terangnya} membuat \omterm{def:b1:water-small-molecules:atp}{ATP} dengan cara ini di
dalam \omterm{def:b1:flowering-plant-organization:organs}{daunnya} (itu menuntut zat pemutus gandengan dan pengukuran pH).
Dengan sebuah zat pemutus gandengan gradiennya runtuh sebelum sintasenya
sanggup memakainya: tak ada \omterm{def:b1:water-small-molecules:atp}{ATP}.
\end{solution}

\begin{solution}{exo:b1:photosynthesis:8}
$3\times 5 + 3\times 1 = 18$ karbon; enam 3-PGA $= 18$; enam G3P $= 18$;
satu G3P keluar (3) dan lima G3P (15) membentuk kembali tiga RuBP (15):
$3 +
15 = 18$. Seimbang.
\end{solution}

\begin{solution}{exo:b1:photosynthesis:9}
Tanpa $\mathrm{CO_2}$ \omterm{def:b1:photosynthesis:fixation}{RuBisCO} berhenti: RuBP-nya tak lagi terpakai
(naik) dan 3-PGA-nya tak lagi dibuat (turun) sementara cahayanya terus
mereduksinya habis. Tanpa cahaya tak ada \omterm{def:b1:water-small-molecules:atp}{ATP} maupun NADPH: 3-PGA-nya
tak lagi direduksi (naik) dan RuBP-nya tak lagi dibentuk kembali
(turun) sementara fiksasinya berlanjut sebentar.
\end{solution}

\begin{solution}{exo:b1:photosynthesis:10}
$E = hc/\lambda = 6.63\times 10^{-34}\times 3\times 10^8/6.8\times
10^{-7} = \qty{2.92e-19}{J}$; tiap mol \qty{176}{kJ}. 48 foton:
\qty{8450}{kJ}; efisiensinya $2870/8450 = \qty{34}{\%}$. Sebuah tanaman
pertanian kehilangan cahaya yang tak terserap (separuh spektrumnya,
pemantulan, penerusan, celah antartumbuhan), fotorespirasinya,
respirasi \omterm{def:b1:flowering-plant-organization:organs}{daun}, \omterm{def:b1:flowering-plant-organization:organs}{batang} dan \omterm{def:b1:flowering-plant-organization:organs}{akarnya}, penjenuhan cahaya \omterm{def:b1:enzymes:enzyme}{enzimnya} pada
matahari penuh, serta musim ketika \omterm{def:b1:flowering-plant-organization:organs}{daunnya} tak ada.
\end{solution}

\begin{solution}{exo:b1:photosynthesis:11}
Tiap 100 karboksilasi, 33 oksigenasi yang melepas $16.7$ karbon: bersih
$83.3$, kehilangan \qty{17}{\%}. Biaya menghindarinya: 2 \omterm{def:b1:water-small-molecules:atp}{ATP} $\approx
0.2$ karbon tiap $\mathrm{CO_2}$, yakni \qty{20}{\%}; rancangan C4
menguntungkan hanya ketika kehilangannya melebihi sekitar seperlima ---
pada keadaan panas dan kering, bukan pada keadaan sejuk.
\end{solution}

\begin{solution}{exo:b1:photosynthesis:12}
Keduanya memakai membran yang memompa proton dengan energi pengangkutan
elektron dan sebuah \omterm{thm:b1:photosynthesis:chemiosmosis}{ATP sintase} yang membelanjakan gradiennya:
mekanismenya sama, dan sintasenya sekerabat. Tetapi elektronnya
berjalan ke arah yang berlawanan --- dari air ke NADPH dalam
\omterm{def:b1:photosynthesis:chloroplast}{kloroplasnya}, diangkat cahaya; dari NADH ke oksigen dalam
\omterm{def:b1:eukaryotic-cell:mitochondrion}{mitokondrianya}, menurun
(\cref{ch:b1:respiration-fermentation}) --- dan daur karbonnya bukan
kebalikan satu sama lain: \omterm{def:b1:photosynthesis:fixation}{daur Calvin} dan daur Krebs tidak berbagi \omterm{def:b1:enzymes:enzyme}{enzim}
mana pun, dan reduksi $\mathrm{CO_2}$ memakai NADPH sedangkan
pelepasannya memakai $\mathrm{NAD^+}$. Persamaan keseluruhannya memang
berkebalikan; mesinnya adalah warisan bersama yang dipakai pada dua
arah, dengan kimia yang berbeda pada ujung karbonnya.
\end{solution}

\begin{solution}{pb:b1:photosynthesis:1}
\textbf{1.} $hc/\lambda = \qty{2.92e-19}{J}$; $\times N_A = \qty{176}{kJ/mol}$.
\textbf{2.} \qty{700}{nm}: \qty{171}{kJ/mol}; \qty{450}{nm}:
\qty{266}{kJ/mol}. Sebuah foton biru mengeksitasi \omterm{def:b1:photosynthesis:pigments}{klorofil} ke keadaan
yang lebih tinggi yang meluruh dalam beberapa pikosekon ke keadaan
tereksitasi terendah yang sama yang dicapai foton merah; kelebihannya
hilang sebagai kalor, dan kimianya melihat satu foton pada kedua hal.
\textbf{3.} $1/2.92\times 10^{-19} = \num{3.4e18}$ foton tiap joule.
\textbf{4.} $\Delta E = 0.82 - (-0.32) = \qty{1.14}{V}$: $\Delta G =
96\,500\times 1.14 = \qty{110}{kJ}$ tiap mol elektron.
\textbf{5.} Dua foton: sekitar $176 + 171 = \qty{347}{kJ}$; tersimpan
\qty{110}{kJ}: \qty{32}{\%}.
\textbf{6.} Empat elektron tiap $\mathrm{O_2}$ (dua molekul air), yang
memberikan dua NADPH: delapan foton.
\textbf{7.} 12 NADPH: 24 elektron, 48 foton.
\textbf{8.} $2.303\times 8.314\times 298\times 3/96\,500 = \qty{0.177}{V}$;
$\qty{17.1}{kJ}$ tiap mol proton.
\textbf{9.} $4 + 8 = 12$ proton tiap $\mathrm{O_2}$: 3 \omterm{def:b1:water-small-molecules:atp}{ATP}.
\textbf{10.} $6\times 3 = 18$ \omterm{def:b1:water-small-molecules:atp}{ATP}: tepat kebutuhan daurnya menurut
perhitungan ini; dalam praktik hasilnya lebih dekat 2.6 tiap
$\mathrm{O_2}$ (\qty{4.7}{H^+} tiap \omterm{def:b1:water-small-molecules:atp}{ATP} dalam \omterm{def:b1:photosynthesis:chloroplast}{kloroplasnya}) dan aliran
siklik mengelilingi \omterm{def:b1:photosynthesis:zscheme}{fotosistem} I memasok sisanya.
\textbf{11.} $4\times 17.1 = \qty{68}{kJ}$ bagi satu \omterm{def:b1:water-small-molecules:atp}{ATP} \qty{50}{kJ}:
\qty{73}{\%}.
\textbf{12.} $10^{-5}\,\mathrm{mol/L}\times 10^{-15}\,\mathrm{L}\times
6\times 10^{23} = 6$ proton bebas dalam lumennya berhadapan dengan
$10^5$ tiap sekon lewat sintasenya: fluksnya dibawa proton yang terikat
pada gugus \omterm{def:b1:water-small-molecules:buffer}{penyangga} \omterm{def:b1:proteins:peptide}{protein} dan \omterm{def:b1:lipids:fattyacid}{lipid} lumennya, yang melepasnya
secepat proton itu pergi.
\textbf{13.} $18\times 50 + 12\times 220 = 900 + 2640 = \qty{3540}{kJ}$
bagi \qty{2870}{kJ} yang tersimpan: \qty{670}{kJ} (seperlima) dilesapkan
sebagai kalor dalam reaksi daurnya, yang harus menurun agar berjalan.
\textbf{14.} $48\times 176 = \qty{8450}{kJ}$; $2870/8450 = \qty{34}{\%}$.
\textbf{15.} $0.34\times 0.45\times 0.9 = \qty{13.8}{\%}$.
\textbf{16.} $400\times 0.005\times 36\,000 = \qty{72}{kJ}$; glukosanya
$0.138\times 72\,000/2\,870\,000 = \qty{3.5e-3}{mol} = \qty{0.62}{g}$.
\textbf{17.} Bersih $0.6\times 0.62 = \qty{0.37}{g}$ glukosa,
$\qty{0.15}{g}$ karbon ($72/180$).
\textbf{18.} Cahaya yang jatuh di antara tumbuhannya atau di tanah;
\omterm{def:b1:flowering-plant-organization:organs}{daun} yang ternaungi \omterm{def:b1:flowering-plant-organization:organs}{daun} lain bekerja di bawah kemampuannya, dan \omterm{def:b1:flowering-plant-organization:organs}{daun}
di bawah matahari penuh yang \omterm{def:b1:lipids:fattyacid}{jenuh} (\omterm{def:b1:enzymes:enzyme}{enzimnya} tak sanggup memakai seluruh
fotonnya); respirasi \omterm{def:b1:flowering-plant-organization:organs}{akar}, \omterm{def:b1:flowering-plant-organization:organs}{batang} dan pada malam hari; fotorespirasi;
minggu-minggu sebelum kanopinya menutup dan sesudah ia menua.
\textbf{19.} 33 oksigenasi, 16.7 karbon hilang, bersih 83.3 diperoleh.
\textbf{20.} Penyelamatannya $33\times 3.5 = 117$ \omterm{def:b1:water-small-molecules:atp}{ATP}; daurnya 300 \omterm{def:b1:water-small-molecules:atp}{ATP};
seluruhnya 417 \omterm{def:b1:water-small-molecules:atp}{ATP} bagi 83.3 karbon: 5.0 \omterm{def:b1:water-small-molecules:atp}{ATP} tiap karbon bersih.
\textbf{21.} $3 + 2 = 5$ \omterm{def:b1:water-small-molecules:atp}{ATP} tiap karbon, tanpa karbon yang hilang ---
setara dalam \omterm{def:b1:water-small-molecules:atp}{ATP}, tetapi tumbuhan C3-nya juga telah kehilangan
\qty{17}{\%} daya karboksilasinya; kira-kira seri pada
\qty{25}{\celsius}.
\textbf{22.} 50 oksigenasi, 25 karbon hilang, bersih 75;
penyelamatannya 175 \omterm{def:b1:water-small-molecules:atp}{ATP}; $475/75 = 6.3$ \omterm{def:b1:water-small-molecules:atp}{ATP} tiap karbon bersih
berhadapan dengan 5 milik \omterm{def:b1:photosynthesis:c4cam}{tumbuhan C4}, dengan seperempat karbonnya
hilang: C4 menang jelas.
\textbf{23.} Melipattigakan $\mathrm{CO_2}$ menggeser persaingan
\omterm{def:b1:photosynthesis:fixation}{RuBisCO} ke arah karboksilasi, sehingga memangkas fotorespirasi dan
menaikkan fiksasi bersih pada tumbuhan C3; \omterm{def:b1:photosynthesis:c4cam}{tumbuhan C4} sudah
menjenuhkan \omterm{def:b1:photosynthesis:fixation}{RuBisCO-nya} dengan $\mathrm{CO_2}$ yang pekat dan tidak
memperoleh apa pun selain sedikit penghematan air.
\textbf{24.} \qty{1}{g} karbon adalah \qty{0.083}{mol}: C3, $250\times 0.083 = \qty{21}{mol}$ air, \qty{375}{g}; CAM, \qty{37}{g}.
\textbf{25.} Sekitar 8 foton tiap $\mathrm{CO_2}$ (48 tiap glukosa);
\qty{34}{\%} dari foton yang terserap sampai glukosa; sekitar
\qty{14}{\%} dari sinar matahari datang sampai gula kasarnya dan
\qty{8}{\%} sampai gula bersih \omterm{def:b1:flowering-plant-organization:organs}{daunnya} --- sebuah angka yang diturunkan
kehilangan seluruh tumbuhan sepanjang seluruh musim menjadi \qty{1}{\%}.
\end{solution}
